\def\baselinestretch{1}

\nonumchapter{Symbols and notation} \vspace{-1cm}

\setlength{\tabcolsep}{0.5cm}
\begin{tabular}{rl}
        \qed & Fin de una prueba.\\
        $\dot{x},\dot{y}$ & El punto es utilizado para denotar la derivada con respecto del tiempo,\\&ejemplo $\dot{x}=x^{(1)}=dx/dt$\\
        $\hat{x},\hat{z}$ & El circunflejo es utilizado para denotar la estimacion de una\\&variable mediante observadores.\\
        $A^T,x^T$ & La transpuesta de la matriz $A$ y el vector $x$.\\
        $\Real$ & El campo de los numeros reales.\\
        $\Complex$ & El campo de los numeros complejos.\\
        $\infty$ & Infinito.\\
        $rango(A)$ & El rango de la matriz A.\\
        $\triangleq$ & Igual por definicion.\\
        $\approx$ & Aproximadamente igual.\\
        $\equiv$ & Identico.\\
        $dist\{x,y\}$ & La distancia entre los puntos $x$ y $y$, en algun espacio definido.\\
        $\inf x$ & El \textit{infimum} de $x$.\\
        $\sup x$ & El \textit{supremum} de $x$.\\
        $\norm{x}$ & Una norma de $x$.\\
        $\lim_{a\rightarrow b}x$ & El limite de $x$ cuando $a$ tiende a $b$.\\
        $\ln$ & Logaritmo natural (en base $e$) de $x$.\\
        $\log x$ & Logaritmo en base 10 de $x$.\\
        $\dim X$ & La dimension del conjunto $X$.\\
        $\mathcal{C}^{\infty}$ & El conjunto de la funciones continuamente diferenciables.\\
        $A^{-1}$ & Inversa de A.\\
        $j$ & Unidad de los numeros imaginarios, $j=\sqrt{-1}$.\\
        $\partial F/\partial x$ & Derivada parcial de la funcion $F$ con respecto a la variable $x$.\\
        $L_fh$ & Derivada Lie de $h$ a lo largo del campo vectorial $f$.\\
        \textit{Cursivas} & El uso de cursivas sirve para denotar terminos en un idioma diferente al Espanol.\\
        \textbf{Negritas} & El uso de negritas sirve para denotar terminos tecnicos.
\end{tabular}
